\begin{abstract}
Text detoxification requires models to reduce harmful or insulting language while
preserving meaning, fluency, and the input language. The problem is especially
difficult in multilingual settings, where the internal representations supporting
rewriting may be partly shared across languages but unevenly transferable. We study
toxic-to-detox representations and activation steering in a nine-language corpus using
mT5, mT0, Llama-3-8B, Qwen2.5-7B-Instruct, and Aya-23-8B.

Our primary tracing analyses use multi-seed probes on Aya and Qwen. At layer 20, Aya
reaches mean pooled toxic-versus-detox AUC of 0.8382 across two fine-tuning seeds and
Qwen reaches 0.8427 across three seeds. Transfer from seven higher-resource languages
to isiXhosa and Yoruba is consistently lower: 0.5966 for Aya and 0.6054 for Qwen.
This gap shows that a representation signal shared across the corpus does not transfer
uniformly to the two lower-resource targets. We also analyze sparse autoencoder
candidate features and activation-steering artifacts. In an exploratory 50-example
Amharic mT5 sweep, steering strength produces a non-monotonic
reference-overlap profile while outputs remain non-empty and unique.

We separate representation decodability, intervention-induced output change, and
validated detoxification improvement. This separation motivates a controlled
train--development--test steering protocol and blinded human evaluation of
non-toxicity, meaning preservation, fluency, and acceptability. Human-evaluation and
controlled-confirmation results will be inserted in the final version.
\end{abstract}
